\chapter{Introduction}
\label{ch:intro}

A switching converter draws its input current in pulses. Left alone, those
pulses flow back out onto the supply bus as conducted noise. An input filter
sits between the bus and the converter to keep that noise in, and it has two
jobs that pull against each other:

\begin{enumerate}
	\item \textbf{Emissions.} Attenuate the converter's noise current enough to
	      meet a conducted emissions limit at the switching frequency.
	\item \textbf{Stability.} Do not destabilize the converter. A lightly damped
	      LC filter in front of a regulated converter can oscillate.
\end{enumerate}

These notes build up the analysis one piece at a time: one LC stage,
damping, two stages, the effect of the source, the common-mode path, and
turn-on. Each chapter applies its results to the same example design.

\section{The example design}
\label{sec:example-spec}

Throughout, the device is a DC--DC converter on a \qty{120}{\volt} bus.

\begin{center}
	\begin{tabular}{@{}ll@{}}
		\toprule
		Quantity                              & Value                                                            \\
		\midrule
		Bus, nominal                          & \qty{120}{\volt} DC                                              \\
		Minimum at converter terminals        & \qty{100}{\volt} DC                                              \\
		Converter output power                & \qty{100}{\watt}                                                  \\
		Efficiency                            & 0.85, so $P_{in} = \qty{117.6}{\watt}$                           \\
		Lowest switching frequency at load    & \qty{400}{\kilo\hertz}                                           \\
		DM emission limit at \qty{400}{\kHz}  & 50~\dBuA{} (example value)                                       \\
		Margin to the limit                   & 6~dB                                                             \\
		Stability target                      & $\abs{\Zo}$ at least 10~dB below $\abs{\Zin}$ at every frequency \\
		Converter overvoltage shutdown        & \qty{200}{\volt}                                                 \\
		\bottomrule
	\end{tabular}
\end{center}

All capacitor values in the examples are \emph{effective} values at 100 to
\qty{120}{\volt} bias. Ceramic capacitors can lose half their capacitance or
more at rated bias, so derate them before you start or every answer shifts.

\section{Notation}

\begin{itemize}
	\item $s = j\omega$, $\omega = 2\pi f$.
	\item $f_0 = 1/(2\pi\sqrt{LC})$ is the resonant frequency of an $L$--$C$ pair.
	\item $R_0 = \sqrt{L/C}$ is its characteristic impedance.
	\item $Q = R_0/\Rloop$, where $\Rloop$ is the series loop resistance.
	\item $\parallel$ means parallel combination: $Z_1 \parallel Z_2 = Z_1Z_2/(Z_1+Z_2)$.
\end{itemize}

$R_0$ and $\Zo$ look alike but are different things. $R_0$ (zero) is a single
number belonging to one $L$--$C$ pair. $\Zo$ (letter o) is a function of
frequency for the whole filter. \zcref{ch:quickref} has the full symbol table.

% -------------------------------------------------------------------

\section{One ladder, two functions}
\label{sec:method}

\begin{idea}
	Every input filter question reduces to two functions of one ladder network:
	\begin{itemize}
		\item the noise transfer $H(s) = i_{line}/i_n$, for emissions, and
		\item the output impedance $\Zo(s)$ seen by the converter, for stability.
	\end{itemize}
	Both come from the same impedance algebra and share the same poles, so a
	resonance that hurts one hurts the other.
\end{idea}

\zcref[S]{fig:dm-circuit} is the two-stage filter used in most of these notes.
Build the impedances left to right, from the source. Divide the noise current
right to left, from the converter.

\begin{figure}[ht]
	\centering
	\begin{circuitikz}[scale=0.85, transform shape]
		\draw (0,0) to[V, v<=$V_s$] (0,3)
		to[generic, l=$Z_s$, i>^=$i_{line}$] (2.2,3)
		to[L, l=$L_{lk}$] (4,3)
		to[L, l_=$L_1$, *-*] (7,3)
		to[short] (7.5,3)
		to[L, l=$L_2$, *-] (10,3)
		to[short, -o] (12,3);
		\draw (4,3) -- (4,4.4) to[R, l=$R_d$] (5.5,4.4) to[L, l=$L_d$] (7,4.4) -- (7,3);
		\draw (7.5,3) to[C, l=$C_1$] (7.5,0);
		\draw (10,3) to[C, l_=$C_2$, *-*] (10,0);
		\draw (0,0) -- (12,0);
		\draw[dashed] (12.2,3.3) rectangle (14.2,-0.3);
		\draw (12,3) -- (13.2,3) to[I, l=$i_n$] (13.2,0) -- (12,0);
		\node[align=center, font=\small] at (13.2,3.7) {converter};
		\draw[-latex, thick] (11.8,1.5) -- node[above, font=\small] {$\Zo$} (10.9,1.5);
		\node[font=\small, align=center] at (2.9,4.6) {CM choke\\(DM leakage)};
	\end{circuitikz}
	\caption{Differential-mode (DM) equivalent of a two-stage input filter.
		Parasitic DCR and ESR are folded into each $Z$.}
	\label{fig:dm-circuit}
\end{figure}

\subsection{The DM equivalent circuit}

\begin{enumerate}
	\item Fold the return line into the loop. Every series element counts once
	      per loop: both DCRs, the choke's total leakage, any series capacitors.
	\item A common-mode (CM) choke contributes only its DM leakage $L_{lk}$ and
	      winding resistance. Its CM inductance cancels for DM current (\zcref{sec:cm-winding}).
	\item The source is an impedance $Z_s$. A stiff bus is a small $R$ plus $L$.
	      Test setups add their own impedance (\zcref{ch:source}).
	\item The converter gets two models. For emissions it is a current source
	      $i_n$: it forces its pulsed input current whatever the filter does.
	      For stability it is a negative resistance.
\end{enumerate}

\begin{derivation}
	A regulated converter holds its output power constant, so below its loop
	crossover it holds $P_{in}$ constant too:
	\[
		I = \frac{P_{in}}{V}
		\quad\Rightarrow\quad
		\frac{dI}{dV} = -\frac{P_{in}}{V^2}
		\quad\Rightarrow\quad
		\Zin = \frac{dV}{dI} = -\frac{V^2}{P_{in}}.
	\]
	Raise the input voltage and the converter draws \emph{less} current.
	That is a negative incremental resistance, and a negative resistance can
	cancel the losses that damp a filter resonance.
\end{derivation}

Above the converter's crossover $\abs{\Zin}$ rises, so using $V^2/P_{in}$ at
every frequency is conservative. The worst case is the smallest $\abs{\Zin}$:
minimum input voltage, maximum load.

\subsection{Building the ladder}

From the source end:
\begin{equation}\label{eq:ladder}
	\begin{aligned}
		Z_{L1}'   & = Z_{L1} \parallel (R_d + sL_d),  & Z_{up,0} & = Z_s + Z_{cm} + Z_{L1}',                   \\
		Z_{up,1}  & = Z_{C1} \parallel Z_{up,0},      & \Zo      & = Z_{C2} \parallel \left(Z_{L2} + Z_{up,1}\right).
	\end{aligned}
\end{equation}

Each $Z$ includes its parasitic:
\[
	Z_L = R_{\DCR} + sL, \qquad Z_C = \ESR + \frac{1}{sC}.
\]

\subsection{Noise transfer}

When a current $i_{in}$ splits between branches $Z_1$ and $Z_2$, the share in
branch 1 is
\begin{equation}\label{eq:current-divider}
	i_1 = i_{in}\,\frac{Z_2}{Z_1 + Z_2}.
\end{equation}

The noise current splits at each shunt capacitor. Multiply one divider per
stage, working from the converter toward the source:
\begin{equation}\label{eq:H-ladder}
	H(s) = \frac{i_{line}}{i_n}
	= \underbrace{\frac{Z_{C2}}{Z_{C2} + Z_{L2} + Z_{up,1}}}_{\text{at } C_2}
	\cdot \underbrace{\frac{Z_{C1}}{Z_{C1} + Z_{up,0}}}_{\text{at } C_1}.
\end{equation}

The emission at a frequency, in \dBuA, is the converter's noise level plus
$20\log_{10}\abs{H}$.

\begin{idea}
	By reciprocity, the same $H$ is the bus-voltage to converter-voltage transfer
	with the converter port open. A peak in $H$ therefore also amplifies bus
	transients and anything injected on the line.
\end{idea}

\subsection{Middlebrook's stability criterion}

\begin{derivation}
	Replace the filter by its Th\'evenin equivalent as seen from the converter:
	a source $v_{th}$ behind $\Zo$. The converter is the impedance $\Zin$.
	The voltage the converter actually gets is
	\[
		v_{in} = v_{th}\,\frac{\Zin}{\Zin + \Zo} = \frac{v_{th}}{1 + \Zo/\Zin}.
	\]
	This has the same form as a closed-loop transfer function $1/(1+T)$, with
	the \emph{minor loop gain}
	\[
		T_m = \frac{\Zo}{\Zin}.
	\]
	The filter-plus-converter system is unstable if the Nyquist plot of $T_m$
	encircles $-1$. If $\abs{T_m} < 1$ at every frequency, the plot stays inside
	the unit circle and cannot reach $-1$. So $\abs{\Zo} < \abs{\Zin}$
	everywhere guarantees the filter cannot destabilize the converter.
\end{derivation}

Design for margin at minimum input voltage and maximum load, with every source
impedance the device will see:
\begin{equation}\label{eq:middlebrook}
	\abs{\Zin}_{min} = \frac{V_{min}^2}{P_{in}}, \qquad
	\text{margin (dB)} = 20\log_{10}\frac{\abs{\Zin}_{min}}{\max_f \abs{\Zo(f)}}.
\end{equation}

\begin{examplebox}[label=ex:zin]{stability target}
	For the example device at minimum input voltage and full load:
	\[
		P_{in} = \frac{\qty{100}{\watt}}{0.85} = \qty{117.6}{\watt}, \qquad
		\abs{\Zin}_{min} = \frac{(\qty{100}{\volt})^2}{\qty{117.6}{\watt}} = \qty{85.0}{\ohm}.
	\]
	A 10~dB margin is a factor of $10^{10/20} = 3.16$, so the filter must keep
	\[
		\abs{\Zo} \leq \frac{\qty{85.0}{\ohm}}{3.16} = \qty{26.9}{\ohm}
	\]
	at every frequency, on every source.
\end{examplebox}

\subsection{Paper shortcuts}

Each is derived in the chapter noted.

\begin{enumerate}
	\item One stage, above $f_0$: $\abs{H} \approx (f_0/f)^2$, a slope of
	      $-40$~dB/decade (\zcref{sec:above-resonance}).
	\item Two stages, above both resonances: $\abs{H} \approx (f_{01}f_{02}/f^2)^2$,
	      where $f_{0k} = 1/(2\pi\sqrt{L_kC_k})$ are the isolated stage resonances.
	      This holds however strongly the stages interact (\zcref{sec:hf-asymptote}).
	\item A lightly damped series-loop resonance: $\Zo$ peak $\approx QR_0 = R_0^2/\Rloop$
	      (\zcref{sec:zo-peak}).
	\item $C_2$ resonating against $L_2$ in series with $C_1$: $\Zo$ peak
	      $\approx X_{C2}^2/\Rloop$, with $X_{C2} = 1/(\omega C_2)$ at that
	      resonance (\zcref{sec:undamped-stage}).
	\item An $R_d$--$L_d$ leg across $L_1$ leaves $L_1 \parallel L_d = L_1 n/(1+n)$
	      at high frequency. It costs $20\log_{10}((1+n)/n)$~dB of attenuation
	      (\zcref{sec:rd-ld}).
\end{enumerate}
